\chapter{Implementation}

The design phase established what the ICMS platform needed to do and how its components should relate to each other. This chapter documents the process of actually building it. Translating the architectural diagrams from Chapter~4 into working code involved some practical realities that shaped our final decisions: hardware constraints, Docker networking quirks, and the challenge of wiring together a Python backend, a JavaScript frontend, and an external SIEM running on a separate physical machine. The following sections walk through each layer of the implementation, grounded in the actual codebase.

All development was done on consumer-grade hardware running Windows and Ubuntu in a distributed home-lab setup. The Wazuh SIEM manager ran on a dedicated laptop at a fixed IP address (\url{192.168.100.100}), and the full backend, frontend, and database stack was containerized on a second machine using Docker Compose. Every service configuration, environment variable, and migration step described here corresponds directly to the files in the project repository.

\section{System Architecture}

As established in the design phase (as shown below in Figure \ref{fig:tier_architecture_impl}), the ICMS platform follows an N-tier client-server architecture.

\begin{figure}[H]
    \centering
    \resizebox{0.85\textwidth}{!}{
    \begin{tikzpicture}[
        tier/.style={rectangle, draw=blue!80!black, thick, fill=blue!5, minimum width=8.5cm, minimum height=1.6cm, align=center, font=\bfseries},
        arrow/.style={-Stealth, thick, draw=black!80},
        label/.style={font=\scriptsize, inner sep=2pt}
    ]
        % Tiers with increased separation to resolve congestion
        \node[tier] (client) at (0, 8.5) {Client \& Presentation Tier\\(React SPA)};
        \node[tier] (app) at (0, 5.0) {Application \& Middleware Tier\\(Django REST Framework)};
        
        \node[tier, minimum width=4cm] (async) at (-3.2, 1.5) {Async Processing\\(Celery \& Redis)};
        \node[tier, minimum width=4cm] (data) at (3.2, 1.5) {Data Persistence\\(PostgreSQL)};
        
        \node[tier, fill=gray!10, draw=gray!80!black] (wazuh) at (-3.2, -1.5) {External Security\\(Wazuh SIEM)};

        % Client <-> App Connections
        \draw[arrow] ([xshift=-2cm]client.south) -- ([xshift=-2cm]app.north) 
            node[midway, left=8pt, label] {HTTPS / JSON};
        \draw[arrow] ([xshift=2cm]app.north) -- ([xshift=2cm]client.south) 
            node[midway, right=8pt, label] {API Responses};

        % App -> Async / Data Connections (Widened entry points)
        \draw[arrow] ([xshift=-2cm]app.south) -- ([xshift=0.8cm]async.north) 
            node[midway, left=10pt, label] {Enqueue Tasks};
        \draw[arrow] ([xshift=2cm]app.south) -- ([xshift=-0.8cm]data.north) 
            node[midway, right=10pt, label] {ORM Queries};

        % Async -> Wazuh (Increased vertical gap)
        \draw[arrow] (async.south) -- (wazuh.north) 
            node[midway, left=8pt, label] {REST API Polling};
        
        % Async -> Data (Increased horizontal gap)
        \draw[arrow] (async.east) -- (data.west) 
            node[midway, above=6pt, label] {State Commits};
    \end{tikzpicture}
    }
    \caption{Implemented N-Tier System Architecture}
    \label{fig:tier_architecture_impl}
\end{figure}

In practice, this mapped onto two concrete code repositories: a \url{backend/} directory containing the Django application and all its sub-apps, and a \url{frontend/} directory containing the React SPA. Both are wired together through a \url{docker-compose.yml} file at the project root that defines seven services and a shared bridge network.

The full service topology (as shown below in Figure \ref{fig:deployment_diagram_impl}) illustrates the container environment and network isolation.

\begin{figure}[H]
    \centering
    \resizebox{0.95\textwidth}{!}{
    \begin{tikzpicture}[
        device/.style={
            rectangle, draw=black!90, thick, fill=white,
            minimum width=4cm, minimum height=2.5cm,
            append after command={
                \pgfextra{
                    \filldraw[black!90, thick, fill=gray!15] (\tikzlastnode.north west) -- ++(0.3, 0.4) -- ($(\tikzlastnode.north east)+(0.3,0.4)$) -- (\tikzlastnode.north east) -- cycle;
                    \filldraw[black!90, thick, fill=gray!15] (\tikzlastnode.north east) -- ($(\tikzlastnode.north east)+(0.3,0.4)$) -- ($(\tikzlastnode.south east)+(0.3,0.4)$) -- (\tikzlastnode.south east) -- cycle;
                }
            }
        },
        artifact/.style={
            rectangle, draw=black!80, thick, fill=blue!5,
            minimum width=3.4cm, minimum height=1cm,
            align=center, font=\scriptsize\bfseries, rounded corners=2pt
        },
        conn/.style={
            thick, draw=black!80, Stealth-Stealth
        },
        label/.style={
            font=\scriptsize, fill=white, inner sep=2pt
        }
    ]
        % ----------------------------------------------------
        % CENTER: Docker Host Server
        % ----------------------------------------------------
        \node[device, minimum width=9cm, minimum height=6cm] (server) at (0, 0) {};
        \node[anchor=north, font=\small\bfseries, yshift=-2mm, align=center] at (server.north) {\textless\textless device\textgreater\textgreater\\Docker Host Server (Ubuntu)};
        
        \node[artifact] (django) at (-2.4, 1) {\textless\textless container\textgreater\textgreater\\Django API Gateway};
        \node[artifact] (celery) at (2.4, 1) {\textless\textless container\textgreater\textgreater\\Celery Async Worker};
        \node[artifact] (postgres) at (-2.4, -1.5) {\textless\textless container\textgreater\textgreater\\PostgreSQL Database};
        \node[artifact] (redis) at (2.4, -1.5) {\textless\textless container\textgreater\textgreater\\Redis Message Broker};

        % Internal Bridge Network Connections
        \draw[dashed, -Stealth, thick, draw=black!70] (django.south) -- (postgres.north) node[midway, fill=white, font=\tiny] {TCP 5432};
        \draw[dashed, -Stealth, thick, draw=black!70] (celery.south) -- (redis.north) node[midway, fill=white, font=\tiny] {TCP 6379};
        \draw[dashed, Stealth-Stealth, thick, draw=black!70] (django.east) -- (celery.west) node[midway, below=3pt, fill=white, font=\tiny] {Task IPC};

        % ----------------------------------------------------
        % LEFT: Client Workstation
        % ----------------------------------------------------
        \node[device, minimum width=4cm, minimum height=3cm] (client) at (-8.5, 1.5) {};
        \node[anchor=north, font=\small\bfseries, yshift=-2mm, align=center] at (client.north) {\textless\textless device\textgreater\textgreater\\Client Workstation};
        \node[artifact] (react) at (-8.5, 1) {\textless\textless artifact\textgreater\textgreater\\React SPA Browser};

        % ----------------------------------------------------
        % RIGHT: Wazuh SIEM Node
        % ----------------------------------------------------
        \node[device, minimum width=4cm, minimum height=3cm] (wazuh) at (8.5, 1.5) {};
        \node[anchor=north, font=\small\bfseries, yshift=-2mm, align=center] at (wazuh.north) {\textless\textless device\textgreater\textgreater\\Wazuh SIEM Node};
        \node[artifact] (sca) at (8.5, 1) {\textless\textless artifact\textgreater\textgreater\\SCA Telemetry API};

        % ----------------------------------------------------
        % EXTERNAL NETWORK CONNECTIONS (Perfectly Horizontal at Y=1)
        % ----------------------------------------------------
        \draw[conn] (react.east) -- (django.west) node[midway, below=4pt, label] {HTTPS / TLS 1.3};
        \draw[conn, -Stealth] (celery.east) -- (sca.west) node[midway, below=4pt, label] {REST API Polling};

    \end{tikzpicture}
    }
    \caption{Implemented Service Topology and Container Networking}
    \label{fig:deployment_diagram_impl}
\end{figure}

When the stack boots, Docker brings up the database and Redis first, waits for their health checks to pass, and only then starts the Django backend. This startup ordering was not optional; without it, the backend would crash on boot trying to run migrations against a database that was not yet accepting connections. The Celery worker and Celery Beat scheduler also depend on the backend image and start only after Redis is confirmed healthy.

\subsection{Presentation Layer Implementation (Frontend)}

The frontend is a React 19 Single Page Application built with Vite 8 as the build and development server. Vite was chosen over the older Create React App toolchain because it uses native ES modules during development, which cuts cold-start time from several seconds down to under 200ms. Routing is handled by React Router v7, which uses a declarative \texttt{<Route>} structure to enforce authentication boundaries at the component level.

Every HTTP request to the Django backend goes through an Axios instance configured with request and response interceptors. The request interceptor reads the JWT access token from \texttt{localStorage} and appends it as a \texttt{Bearer} header automatically. The response interceptor watches for \texttt{401 Unauthorized} responses; when one arrives, it triggers a token refresh and, if the refresh also fails, clears the session and redirects to the login page. This means we never had to write token-refresh logic in individual components.

The UI itself is styled with Tailwind CSS v4 using a custom \texttt{@theme} token set for brand colours, spacing, and typography. Data visualisation is handled by Recharts, which renders the compliance score gauge and per-framework bar charts on the main dashboard. Icon assets come from Lucide React. The routing table enforces RBAC at the frontend level too: the \texttt{/manage-admins} route requires the \texttt{required\allowbreak Role=``super\_admin''} prop on its \texttt{ProtectedRoute} wrapper, so Django's RBAC check at the API layer has a corresponding gate on the UI side.

\subsection{Middleware and API Gateway (Backend)}

The backend is a Django 5.1 project with a single installed application, \texttt{compliance}, that contains every model, view, serializer, and management command the platform needs. The project configuration lives in \texttt{core/settings.py}, and all secrets are loaded at startup via \texttt{python-decouple} from a \texttt{.env} file, so no credentials are ever hardcoded in source control.

Authentication is handled by \url{django-restframework-simplejwt}. Access tokens are valid for 24 hours; refresh tokens last seven days with rotation enabled, so each refresh call issues a new refresh token and invalidates the old one. The \texttt{Custom\allowbreak Token\allowbreak Obtain\allowbreak Pair\allowbreak View} overrides the standard login endpoint to add a two-factor authentication gate: if the user has TOTP enabled via \texttt{pyotp}, the view returns an intermediate \texttt{requires\_\allowbreak 2fa} flag instead of issuing tokens, forcing the client to complete a second verification round-trip before credentials are granted.

The \texttt{compliance} app's \url{views.py} contains both DRF ViewSets (for the standard CRUD resources like frameworks, controls, and policies) and custom \texttt{APIView} subclasses for complex operations such as dashboard aggregation and PDF report generation. The \texttt{Dashboard\allowbreak Summary\allowbreak View} is the most query-intensive endpoint in the system. It uses a single annotated queryset with \texttt{DISTINCT ON agent\_id} to find the latest scan per agent, then computes per-framework pass/fail counts using Django ORM \texttt{Count()} with \texttt{Q()} filter objects. This collapses what would otherwise be a na\"{i}ve N+1 loop across twenty-plus framework controls into a single SQL statement.

\subsection{Asynchronous Processing Engine}

Any time a compliance scan runs, it involves several seconds of HTTP polling against the external Wazuh REST API. Running that on the main Django request thread would block the server until the scan finished, effectively taking the dashboard offline for everyone else. The solution is to dispatch the work to a Celery worker process running in a separate container.

The Celery application is initialised in \url{core/celery.py}. It reads its configuration from Django settings using the \texttt{CELERY\_} namespace and calls \texttt{autodiscover\_tasks()} so it automatically picks up any \url{tasks.py} file across registered apps. The only task in the system is \texttt{run\_scheduled\_\allowbreak wazuh\_sync} in \url{compliance/tasks.py}, decorated with \texttt{@shared\_task(bind=True, max\_retries=3)}. The \texttt{bind=True} argument gives the task access to \texttt{self}, which it uses to call \texttt{self.retry(exc=exc, countdown=60)} when a Wazuh API call fails. This means a temporary network hiccup triggers an automatic 60-second backoff and retry rather than a permanent failure.

Redis runs as a separate container on port 6379 and acts as both the message broker and the result backend. The Celery Beat scheduler, also running as its own container, uses the \texttt{Database\allowbreak Scheduler} from \url{django-celery-beat}. This means scan frequency (daily, weekly, or monthly) is stored in the PostgreSQL database and can be changed at runtime through the Settings page in the UI, without restarting any containers. When a manual scan is triggered from the frontend, the \url{/api/run-scan/} endpoint receives a \texttt{POST} request, calls \texttt{run\_scheduled\_\allowbreak wazuh\_sync.delay()}, and immediately returns a \texttt{202 Accepted} response. The scan then runs in the worker container with no impact on the API's availability.

\subsection{Data Persistence Layer}

PostgreSQL 18 is the primary database, running in its own container and accessible to the Django backend and Celery worker over the Docker bridge network on port 5432. The database is never exposed to the host machine. Django connects via \url{psycopg[binary]} version 3.2, which uses the binary distribution of the C-based libpq driver. There is no additional connection pooling middleware; Django's built-in persistent connection support (\texttt{CONN\_\allowbreak MAX\_\allowbreak AGE}) handles connection reuse.

The schema was built entirely through Django ORM migrations. During the first boot of the backend container, the entrypoint script runs \texttt{python manage.py migrate} before starting the WSGI server. This ensures the schema is always in sync with the models before the application accepts traffic. The initial seed data (ISO 27001 controls, SOC 2 controls, and the Wazuh rule mappings between them) is loaded via a custom management command, \texttt{seed\_\allowbreak enterprise\_\allowbreak matrix}, commissioned to read from a structured data file and bulk-create the relational records.

The most performance-sensitive schema decision was storing raw Wazuh SCA check data in a \texttt{JSONField} on the \texttt{ScanResult} model. Rather than trying to normalise the variable-structure SCA response (which can contain arbitrary compliance keys and values) into relational columns, we kept the raw JSON payload intact. This made ingestion logic simpler and preserved the original evidence for audit purposes, while the structured compliance evaluation (pass/fail, score, framework mapping) was stored in the relational columns alongside it.

\subsection{Architectural Technology Stack}

Table~\ref{tab:tech_stack} summarises the full technology stack across all four implementation tiers, mapping each component to its specific library or service and its role in the system.

\begin{table}[H]
    \centering
    \caption{Architectural Technology Stack and Implementation Roles}
    \label{tab:tech_stack}
    \renewcommand{\arraystretch}{1.5}
    \resizebox{\textwidth}{!}{
    \begin{tabular}{|>{\raggedright\arraybackslash}p{3cm}|>{\raggedright\arraybackslash}p{3.5cm}|>{\centering\arraybackslash}p{2cm}|>{\raggedright\arraybackslash}p{7cm}|}
        \hline
        \rowcolor{blue!10}
        \textbf{Architecture Tier} & \textbf{Technology} & \textbf{Version} & \textbf{Implementation Role} \\ \hline
        Client (Frontend) & React & 19.0.0 & Single Page Application (SPA) UI layer. Manages client-side routing and state without full page reloads. \\ \hline
        Client (Frontend) & Vite & 8.0 & Frontend build tool and development server, replacing Webpack for faster Hot Module Replacement (HMR). \\ \hline
        Client (Frontend) & Recharts & 3.8 & Data visualization library used to render the compliance metrics and endpoint status dashboards. \\ \hline
        Middleware (Backend) & Django & 5.1 & Core Python framework managing the ORM, database migrations, and overarching project structure. \\ \hline
        Middleware (Backend) & Django REST Framework & 3.15 & Serializes Django models into JSON payloads and handles ViewSet routing for the React frontend. \\ \hline
        Middleware (Backend) & SimpleJWT & 5.3 & Handles stateless authentication, issuing 24-hour access tokens and 7-day rotating refresh tokens. \\ \hline
        Middleware (Backend) & xhtml2pdf & 0.2.16 & Generates the physical PDF audit reports directly from Django HTML templates on the server side. \\ \hline
        Async Processing & Celery & 5.4 & Worker nodes that handle Wazuh API polling and email dispatching so the main API thread does not block. \\ \hline
        Async Processing & Redis & 7-alpine & In-memory message broker that queues the Celery tasks and temporarily stores the task results. \\ \hline
        Data Persistence & PostgreSQL & 18 & Primary ACID-compliant relational database. Stores the normalized RBAC models and rule mappings. \\ \hline
        Data Persistence & psycopg[binary] & 3.2 & The C-based adapter that allows the Django backend to communicate synchronously with PostgreSQL. \\ \hline
    \end{tabular}
    }
\end{table}

\section{System Internal Components}

The internal architecture of the Intelligent Compliance Management System is designed as a collection of specialized modules that work in concert to manage the entire compliance lifecycle. This modularity ensures a clear separation of concerns, where identity management, data governance, security telemetry ingestion, and reporting are handled by independent logical units. By partitioning the system into these functional modules, we established a framework that is both resilient to external API failures and scalable across distributed infrastructure. The following subsections detail the architectural intent and implementation logic of each primary system component.

\subsection{Identity and Access Management (IAM) Module}

The \texttt{IAM} module serves as the security foundation of the ICMS, enforcing strict authentication and authorization policies across the entire stack. We implemented a custom identity model that extends the core Django user system through a \texttt{UserProfile} entity, providing the necessary granularity for Role-Based Access Control (RBAC). The module recognizes three distinct operational roles: Superadmins, who possess full system and configuration access; Admins, who manage frameworks and oversee compliance activities; and Auditors, who are restricted to read-only views of scan results and report generation.

Authentication is handled through a stateless JSON Web Token (JWT) architecture, which eliminates the need for server-side session storage and improves scalability. Upon successful credential verification, the system issues a short-lived access token and a long-lived rotating refresh token. For accounts requiring elevated security, we integrated a two-factor authentication (2FA) layer based on the Time-based One-Time Password (TOTP) algorithm. The IAM module orchestrates a two-step login sequence where the user must provide a valid 2FA code before the backend issues the final JWT session. This security policy is enforced at the API layer through Django REST Framework permissions and mirrored at the presentation layer using React route guards, ensuring that unauthorized users cannot bypass the interface to access sensitive data.

\subsection{Governance and Compliance Data Engine}

The \texttt{Governance and Compliance Data Engine} manages the complex relational mappings required to translate abstract regulatory requirements into measurable technical controls. At the top level, the engine stores compliance standards, such as ISO 27001 or SOC 2, as \texttt{Framework} records. Each framework is decomposed into a hierarchy of \texttt{Control} entities, which represent the specific security requirements that an organization must meet. To bridge the gap between policy and practice, the engine links these controls to organizational \texttt{Policy} documents through a many-to-many relationship, providing a documented trail of how each technical requirement is satisfied by internal governance.

The most critical architectural feature of this engine is the mapping logic that connects technical security telemetry to regulatory clauses. We engineered a \texttt{WazuhMapping} system that stores unique identifiers for Security Configuration Assessment (SCA) rules and ties them to their respective governance controls. This allows the system to automatically categorize raw security events into the appropriate compliance bucket. Furthermore, the engine implements a risk-scoring algorithm that utilizes weights and thresholds defined in the system settings to calculate real-time compliance percentages. This data structure ensures that any change in the security state of a remote endpoint is immediately reflected in the high-level compliance posture of the entire organization.

\subsection{SIEM Telemetry and Synchronization Module}

The \texttt{SIEM Telemetry and Synchronization Module} is the most architecturally complex component of the ICMS, responsible for the deterministic ingestion of security data from the Wazuh SIEM manager. Given the potentially large volume of security logs and the inherent latency of external API communication, we implemented this module as a multi-tiered asynchronous system. A dedicated Celery worker tier handles the execution of synchronization tasks, utilizing Redis as a high-performance message broker to queue requests and prevent the main application thread from blocking.

The synchronization lifecycle begins with a scheduled polling event managed by the Celery Beat service. The module authenticates with the Wazuh REST API using a temporary bearer token and proceeds to enumerate all active security agents. For each agent, the system fetches the latest SCA scan results, which arrive as unstructured JSON payloads. Rather than attempting to force this variable data into rigid relational columns, we engineered the module to ingest the raw logs into a \texttt{JSONField} while simultaneously extracting the pass/fail outcomes. This dual-storage approach preserves the original audit evidence for compliance purposes while providing indexed data for rapid dashboard rendering. To ensure reliability, the module implements an exponential backoff and retry policy, allowing the system to gracefully handle network timeouts or temporary SIEM downtime without losing telemetry data.

\subsection{Automated Audit and Reporting Module}

The \texttt{Automated Audit and Reporting Module} is designed to transform thousands of raw scan results into immutable, professional-grade audit artifacts. When an administrator or auditor requests a compliance report, the module performs an intensive aggregation query across the historical scan data, calculating compliance percentages for every active framework and identifying specific gaps where controls have failed. This aggregation is optimized through database-level annotations to ensure that reports can be generated for large-scale environments with hundreds of endpoints.

The architectural intent behind this module is the creation of audit-ready evidence that meets the requirements of external regulatory bodies. The system utilizes a server-side rendering engine to convert Django-templated HTML into binary PDF files. These documents include framework-level summaries, endpoint-specific status lists, and gap analysis sections that highlight non-compliant controls. By offloading the PDF generation to the background worker tier, we ensured that the user interface remains responsive even when generating comprehensive 100-page audit dossiers. The resulting files are served through a secure streaming interface, preventing the need for temporary storage of sensitive compliance data on the server's local file system.

\subsection{Frontend State and Visualization Layer}

The \texttt{Frontend State and Visualization Layer} provides the interactive interface through which users manage the compliance lifecycle. Built as a React Single Page Application (SPA), this layer utilizes a component-based architecture to provide a seamless user experience. We engineered a global \texttt{AuthContext} provider that serves as the "brain" of the client-side session, managing JWT token persistence, user role state, and automatic token refresh logic through Axios interceptors. This ensures that every outbound API request is transparently authenticated and that the user's session remains active as long as the refresh window allows.

Data visualization is a primary focus of this layer, as compliance officers require high-level insights into organizational security states. We integrated a specialized charting library to render dynamic compliance score gauges, trend lines, and endpoint status distributions. These visualizations are not static; they are driven by real-time data fetched from the backend aggregation views. Furthermore, the frontend implements dynamic component mounting based on the user's RBAC role. An auditor session, for instance, initializes a completely different set of sidebar navigation and dashboard widgets than an administrator session, ensuring that users only interact with the features relevant to their operational scope. This architectural separation between state management and presentation ensures a maintainable and secure frontend codebase.

\section{Tools and Technologies}
The implementation of the ICMS platform required a diverse technology stack, spanning from core backend infrastructure to advanced AI-assisted development environments. Table \ref{tab:tools_tech} summarizes the primary tools and frameworks utilized throughout the project lifecycle.

\begin{table}[H]
    \centering
    \caption{Summary of Tools, Technologies, and Development Stack}
    \label{tab:tools_tech}
    \renewcommand{\arraystretch}{1.4}
    \resizebox{0.95\textwidth}{!}{
    \begin{tabular}{|>{\raggedright\arraybackslash}p{4cm}|>{\raggedright\arraybackslash}p{4.5cm}|>{\raggedright\arraybackslash}p{6.5cm}|}
        \hline
        \rowcolor{blue!10}
        \textbf{Category} & \textbf{Tool / Technology} & \textbf{Purpose in Project} \\ \hline
        Programming Languages & Python 3.12, JavaScript & Core backend API logic and frontend UI interactivity. \\ \hline
        Backend Frameworks & Django 5.1, DRF & Database ORM mapping, migrations, and RESTful API creation. \\ \hline
        Frontend Frameworks & React 19, Tailwind CSS 4.0 & Component-based SPA architecture and utility-first UI styling. \\ \hline
        Database \& Messaging & PostgreSQL 18, Redis 7 & ACID-compliant data persistence and asynchronous task queuing. \\ \hline
        UI/UX Design & Google Stitch & Iterative wireframing, prototyping, and user interface design. \\ \hline
        Development IDE & Antigravity & Primary orchestration platform and integrated coding environment. \\ \hline
        AI Code Assistance & Claude (Opus 4.6, Sonnet 4.6), Gemini (3.1 Pro High, 3 Flash) & Advanced algorithmic modeling, automated code generation, and debugging support. \\ \hline
        Documentation & LaTeX, MiKTeX, TikZ & Academic report typesetting and mathematical/architectural diagramming. \\ \hline
    \end{tabular}
    }
\end{table}

The selection of these technologies was driven by the specific engineering requirements of a security governance platform. We chose the Python/Django stack for the backend due to its maturity in handling complex relational data through an integrated Object-Relational Mapper (ORM) and its robust support for administrative authentication. On the presentation layer, the React library was utilized to build a responsive Single Page Application (SPA). This choice allowed the system to manage complex UI states, such as real-time dashboard updates and interactive compliance gauges, without requiring full page reloads, thereby improving operational efficiency for auditors and administrators.

Prior to the implementation phase, the system underwent an iterative UI/UX design process using Google Stitch. We utilized this platform for rapid prototyping, developing low-fidelity wireframes and high-fidelity mockups to validate user flows before writing any frontend code. This phase was critical for establishing the "Cyber-SaaS" aesthetic, ensuring that the final React implementation remained consistent with modern design standards and that navigation paths for security assessments were intuitive and logically structured.

The development workflow was characterized by a modern, AI-driven approach centered around the Antigravity integrated development environment. Antigravity served as the orchestration hub, coordinating local codebase management with advanced AI modeling capabilities. Throughout the project, we utilized models such as Claude and Gemini for complex pair-programming tasks, specifically for engineering the Wazuh telemetry parser and optimizing the Celery task retry logic. This collaborative environment enabled the rapid debugging of container networking issues and the generation of structured API documentation, significantly reducing the implementation timeline while maintaining a high standard of code quality.

Finally, the academic documentation of the thesis was exclusively managed through LaTeX, MiKTeX, and TikZ. We utilized these tools to ensure that the final report adhered to strict academic typesetting standards and that all architectural diagrams were mathematically precise. The TikZ library, in particular, allowed for the programmatic generation of high-fidelity system topology diagrams, which maintained visual consistency across both the design and implementation chapters. This documentation stack guaranteed that the complex technical relationships within the ICMS platform were communicated with maximum clarity and professional rigor.
